\subsection{滤过模的 Hausdorff 完备化}

令 G 为一个滤过群，其滤过 \((\mathbf{G}_{n})\) 由 G 的正规子群组成；我们已经在第 6 节回顾过，拓扑群 G 的 Hausdorff 完备化 \(\hat{G}\) 可典范地识别为离散群的逆极限 \(\lim_{\substack{\mathbf{G}/\mathbf{G}_{n}\\ }} \overleftarrow{\mathbf{G}}\)，其中离散群为 \(G/G_{n}\)；典范同态 \(i: G \to \hat{G}\) 的像是与 G 相关的 Hausdorff 群（在 \(\hat{G}\) 中处处稠密），其核是闭包 \(\bigcap G_{n}\)，即 G 中 \(\{0\}\) 的闭包。G 的子群的 Hausdorff 完备化 \(\hat{G}_{n}\)（子群为 \(G_{n}\)）可识别为 \(i(\mathbf{G}_{n})\) 在 \(\hat{G}\) 中的闭包（《一般拓扑学》，第二章 §3，第 9 节，命题 18 的推论 1），并且由于 \(G_{n}\) 在 G 中是闭的，


\[
\mathbf {G} _ {n} = \stackrel {- 1} {i} (\hat {\mathbf {G}} _ {n}) = \stackrel {- 1} {i} (\hat {\mathbf {G}} _ {n} \cap i (\mathbf {G})).\tag{18}
\]

此外，\(\hat{G}_n\) 构成 \(\hat{G}\) 中 0 的邻域基本系（《一般拓扑学》，第三章 §3，第 4 节，命题 7），因而是 \(\hat{\mathbf{G}}\) 的正规开子群（《一般拓扑学》，第三章 § 2，第 3 节，命题 8）；\(\hat{\mathbf{G}}\) 上的拓扑由滤过 \((\hat{\mathbf{G}}_n)\) 定义，按定义总是分离的。由于 \(i(\mathbf{G})\) 在 \(\hat{\mathbf{G}}\) 中稠密且 \(\hat{\mathbf{G}}_n\) 是开集，

\[
\hat {\mathbf {G}} = i (\mathbf {G}). \hat {\mathbf {G}} _ {n}\tag{19}
\]

同样有

\[
\hat {\mathbf {G}} _ {n - 1} = i \left(\mathbf {G} _ {n - 1}\right). \hat {\mathbf {G}} _ {n}.\tag{20}
\]

由 (18) 和 (19) 可得，滤过 \((\hat{G}_n)\) 是穷竭的，当且仅当滤过 \((G_n)\) 是穷竭的。

第二同构定理（《代数》，第一章 § 6，第 13 节，定理 6 (d)）以及等式 (18)、(19) 和 (20) 表明，典范同态

\[
\mathrm{G} _ {n - 1} / \mathrm{G} _ {n} \rightarrow \hat {\mathrm{G}} _ {n - 1} / \hat {\mathrm{G}} _ {n}, \quad \mathrm{G} / \mathrm{G} _ {n} \rightarrow \hat {\mathrm{G}} / \hat {\mathrm{G}} _ {n},\tag{21}
\]

都是双射，因而典范同态

\[
\operatorname{gr} (\mathbf {G}) \rightarrow \operatorname{gr} (\hat {\mathbf {G}}).\tag{22}
\]

现在设 A 是一个滤过环，E 是一个滤过 A-模，而 \((\mathbf{A}_{n})\) 和 \((\mathbf{E}_{n})\) 分别是 A 和 E 的滤过；我们假设这些滤过是穷竭的，从而相应的拓扑使 A 成为拓扑环、E 成为拓扑 A-模（第 5 节，命题 3）。于是我们已经定义了（《一般拓扑学》，第三章 § 6，第 5、6 节）拓扑环 \(\hat{A}\) 和拓扑模 \(\hat{E}\)，后者是拓扑 \(\hat{A}\)-模。若 \(i: A \to \hat{A}\) 是典范同态，则 \(i(\mathbf{A}_{m})i(\mathbf{A}_{n}) \subset i(\mathbf{A}_{m+n})\)，因此由 \(\hat{A}\) 中乘法的连续性，

\[
\hat {\mathrm{A}} _ {m} \hat {\mathrm{A}} _ {n} \subset \hat {\mathrm{A}} _ {m + n}\tag{23}
\]

因为 \(\hat{\mathbf{A}}_n\) 是 \(i(\mathbf{A}_n)\) 在 \(\hat{\mathbf{A}}\) 中的闭包。类似地可证明

\[
\hat {\mathrm{A}} _ {m} \hat {\mathrm{E}} _ {n} \subset \hat {\mathrm{E}} _ {m + n};\tag{24}
\]

换言之：

命题 15. 设 A 是一个滤过环，E 是一个滤过 A-模，A 和 E 的相应滤过 \((\mathbf{A}_{n})\)、\((\mathbf{E}_{n})\) 是穷竭的。那么 \((\hat{\mathbf{A}}_{n})\) 是一个与 \(\hat{A}\) 的环结构相容的滤过，而 \((\hat{\mathbf{E}}_{n})\) 是一个与 \(\hat{E}\) 上的模结构相容的滤过，该模结构来自滤过环 \(\hat{A}\)；此外，这些滤过是穷竭的，并分别定义 \(\hat{A}\) 和 \(\hat{E}\) 上的拓扑。最后，分次 Z-模的典范映射 \(\mathrm{gr}(\mathbf{A}) \to \mathrm{gr}(\hat{\mathbf{A}})\) 和 \(\mathrm{gr}(\mathrm{E}) \to \mathrm{gr}(\hat{\mathrm{E}})\) 分别是分次环同构和分次 gr(A)-模同构。

在下文中，对于每个一致空间 X，\(j_{x}\) 表示从 X 到其 Hausdorff 完备化 \(\hat{X}\) 的典范映射，\(\mathbf{X}_{0}=j_{\mathbf{x}}(\mathbf{X})\) 表示 \(\hat{X}\) 的一致子空间，它是与 X 相关的 Hausdorff 空间。回忆一下，X 上的拓扑是 \(j_{x}\) 下 \(X_{0}\) 上拓扑的逆像（《一般拓扑学》，

第二章 § 3，第 7 节，命题 12）。还要回忆，对于每个一致连续映射 \(f\colon\mathbf{X}\to\mathbf{Y},\hat{f}\) 表示从 \(\hat{\mathbf{X}}\) 到 \(\hat{\mathbf{Y}}\) 的一致连续映射，它满足 \(\hat{f}\circ j_{\mathbf{X}}=j_{\mathbf{Y}}\circ f\)（同上，命题 15）；如果 \(\mathbf{X}\) 是 \(\mathbf{Y}\) 的一致子空间且 \(f\) 是典范嵌入，则 \(\hat{\mathbf{X}}\) 可识别为 \(\hat{\mathbf{Y}}\) 的一致子空间，而 \(\hat{f}\) 是 \(\hat{\mathbf{X}}\) 到 \(\hat{\mathbf{Y}}\) 的典范嵌入（同上，第 9 节，命题 18 的推论 1）。

引理 2. 设 \(\mathbf{X} \xrightarrow{f} \mathbf{Y} \xrightarrow{g} \mathbf{Z}\) 是拓扑群的严格态射的正合列（《代数》，第二章 §1，第 4 节，注）。假设 \(\mathbf{X}, \mathbf{Y}, \mathbf{Z}\) 都有 Hausdorff 完备群，且 \(\mathbf{X}, \mathbf{Y}, \mathbf{Z}\) 的单位元各自都有可数的邻域基本系。那么 \(\hat{\mathbf{X}} \xrightarrow{f} \hat{\mathbf{Y}} \xrightarrow{g} \hat{\mathbf{Z}}\) 是严格态射的正合列。

令 \(N_{f}, N_{g}\) 分别为 f 和 g 的核；记

\[
f = f _ {3} \circ f _ {2} \circ f _ {1}
\]

其中 \(f_1\) 是典范映射 \(\mathbf{X} \to \mathbf{X} / \mathbf{N}_f\)，\(f_2\) 是从 \(\mathbf{X} / \mathbf{N}_f\) 到 \(\mathbf{N}_g\) 的同构，而 \(f_3\) 是典范嵌入 \(\mathbf{N}_g \to \mathbf{Y}\)。我们已经知道 \(\hat{f}_2\) 是从 \((\mathbf{X} / \mathbf{N}_f)^{\wedge}\) 到 \(\hat{\mathbf{N}}_g\) 的同构，并且刚才回顾过 \(\hat{f}_3\) 是从 \(\hat{\mathbf{N}}_g\) 到 \(\hat{\mathbf{Y}}\) 的单射严格态射；如果我们证明 \(\hat{f}_1\) 是满射严格态射，就可推出 \(\hat{f}\) 是严格态射（《一般拓扑学》，第三章 § 2，第 8 节，注 2）。令 \(g_1\) 为典范映射 \(\mathbf{Y} \to \mathbf{Y} / \mathbf{N}_g\)；如果我们证明 \(\hat{g}_1\) 是核为 \(\hat{\mathbf{N}}_g\) 的满射严格态射，那么如上所述可见 \(\hat{g}\) 是严格态射，并且序列 \(\hat{\mathbf{X}}\xrightarrow{f}\hat{\mathbf{Y}}\xrightarrow{\hat{g}}\hat{\mathbf{Z}}\) 将是正合的。因此问题归结为证明：若 \(\mathbf{Y}=\mathbf{X}/\mathbf{N}\)（其中 \(\mathbf{N}\) 是 \(\mathbf{X}\) 的正规子群），且 \(f\colon\mathbf{X}\to\mathbf{Y}\) 是典范映射，则 \(\hat{f}\) 是核为 \(\hat{\mathbf{N}}\) 的满射严格态射。

令 \(f_{0}: X_{0} \to Y_{0}\) 为与 \(\hat{f}\) 在 \(X_{0}\) 上重合的映射；由于 \(j_{X}\)（相应地 \(j_{Y}\)）是从 X 到 \(X_{0}\)（相应地从 Y 到 \(Y_{0}\)）的满射严格态射，\(f_{0}\) 是满射严格态射（《一般拓扑学》，第三章 § 2，第 8 节，注 3）。现在 \(X_{0}\) 和 \(Y_{0}\) 是可度量化的（《一般拓扑学》，第九章 § 3，第 1 节，命题 1）；于是根据《一般拓扑学》，第九章 § 3，第 1 节，命题 4 的推论 1 和引理 1，\(\hat{f}_{0} = \hat{f}\) 是满射严格态射，并且其核是 \(\hat{N}_{0}'\)，它是 \(\hat{X}\) 中核 \(N_{0}'\)（即 \(f_{0}\) 的核）的闭包。于是只需证明 \(\hat{N}_{0}' = \hat{N}\)。现在 \(N_{0}'\) 显然包含 \(N_{0} = j_{X}(N)\)；只需证明 \(N_{0}'\) 包含于闭包 \(\bar{N}_{0}\)，该闭包是在 X 中取 \(N_{0}\) 的闭包。现在，

\[
\mathbf {U} = \bar {j} _ {\mathbf {X}} ^ {- 1} (\mathbf {X} _ {0} - \overline {{\mathbf {N}}} _ {0}) = \mathbf {X} - \bar {j} _ {\mathbf {X}} ^ {- 1} (\overline {{\mathbf {N}}} _ {0})
\]

是 X 中的一个开集，并且不与 N 相交；由于 f 是满射严格态射，\(V = f(U)\) 是 Y 中的一个开集，不包含 Y 的单位元 \(e'\)，因而也不与 \(e'\) 的闭包相交；于是 \(j_{Y}(V)\) 不包含 \(Y_{0}\) 的单位元。但是 \(j_{Y}(V) = f_{0}(X_{0} - \overline{N}_{0})\)，所以 \(N_{0}' \subset \overline{N}_{0}\)，这就完成了引理 2 的证明。

命题 16. 设 A 是一个滤过环，\((\mathbf{A}_n)\) 是其滤过，E 是一个 A-模，且 \((\mathbf{E}_n)\) 是由 A 上的滤过导出的 E 上的滤过，由 \(\mathbf{E}_n = \mathbf{A}_n\mathbf{E}\) 组成。假设滤过 \((\mathbf{A}_n)\) 是穷竭的，且模 E 是有限生成的。如果 \(i: \mathbf{E} \to \hat{\mathbf{E}}\) 是典范映射，那么对所有 \(n \in \mathbf{Z}\)，

\[
\hat {\mathrm{E}} _ {n} = \hat {\mathrm{A}} _ {n} \hat {\mathrm{E}} = \hat {\mathrm{A}} _ {n} i (\mathrm{E}) \quad a n d \quad \hat {\mathrm{E}} = \hat {\mathrm{A}}. i (\mathrm{E}).\tag{25}
\]

特别地，\(\hat{\mathbf{E}}\) 是有限生成的 \(\hat{\mathbf{A}}\)-模。

等式 \(A_{n}E = E_{n}\) 蕴含：由于 \(\hat{A}\)-模上的外部运算连续，可得 \(\hat{E}, \hat{A}_{n}\hat{E} \subset \hat{E}_{n}\)，并且显然 \(\hat{A}_{n}\hat{E} \supset \hat{A}_{n}i(E)\)。根据假设，存在满同态 \(u: L \to E\)，其中 \(L = A_{s}^{I}\)，I 是有限集；给 L 赋予乘积滤过，由 \(L_{n} = A_{n}^{I}\) 组成，该滤过在 L 上定义乘积拓扑；于是 \(\hat{L} = \hat{A}_{s}^{I}\) 且 \(\hat{L}_{n} = \hat{A}_{n}^{I}\)（《一般拓扑学》，第二章 § 3，第 9 节，命题 18 的推论 2）。令 \(j: L \to \hat{L}\) 为典范映射，\((e_{i})_{i \in I}\) 为 L 的典范基；对于元素 \(\sum_{i \in I} a_{i}j(e_{i})\)（其中 \(a_{i} \in \hat{A}\) 对所有 \(i \in I\) 都成立）属于 \(\hat{L}_{n}\)，其充要条件是对所有 i 都有 \(a_{i} \in \hat{A}_{n}\)；于是 \(\hat{L}_{n} = \hat{A}_{n}.j(L)\)。既然如此，按定义 \(u(\mathbf{L}_{n}) = \mathbf{A}_{n}\mathbf{E} = \mathbf{E}_{n}\)，所以 u 是从 L 到 E 的严格态射（《一般拓扑学》，第三章 § 2，第 8 节，命题 24）。引理 2 随即表明，\(\hat{u}: \hat{L} \to \hat{E}\) 是满射严格态射。由于 \(\hat{L}_{n}\) 是 \(\hat{L}\) 的开子群，\(\hat{u}(\hat{\mathbf{L}}_{n})\) 是 \(\hat{E}\) 的开（因而闭）子群；但 \(\hat{u}(\hat{\mathbf{L}}_{n}) = \hat{\mathbf{A}}_{n}\hat{u}(j(\mathbf{L})) = \hat{\mathbf{A}}_{n}i(\mathbf{E})\)，并且由于 \(i(\mathbf{E}_{n}) \subset \mathbf{A}_{n}i(\mathbf{E}) \subset \hat{\mathbf{A}}_{n}i(\mathbf{E})\)，最终有 \(\hat{\mathbf{E}}_{n} \subset \hat{\mathbf{A}}_{n}i(\mathbf{E}) \subset \hat{\mathbf{A}}_{n}\hat{\mathbf{E}} \subset \hat{\mathbf{E}}_{n}\)，因而 \(\hat{\mathbf{E}}_{n} = \hat{\mathbf{A}}_{n}\hat{\mathbf{E}} = \hat{\mathbf{A}}_{n}i(\mathbf{E})\)；令 n = 0，得到 (25) 的第二个公式。

推论 1. 在命题 16 的条件下，若 A 是完备的，则 E 也是完备的。

由于此时典范映射 \(\mathbf{A} \to \hat{\mathbf{A}}\) 是满射（第 6 节，命题 5），由 (25) 得 \(\hat{\mathbf{E}} = i(\mathbf{E})\)，结论由第 6 节的命题 5 得出。

推论 2. 设 A 是一个交换环，m 是 A 的有限生成理想，\(\hat{\mathbf{A}}\) 是 A 关于 m-进拓扑的 Hausdorff 完备化。那么 \(\widehat{\mathfrak{m}^n} = (\hat{\mathfrak{m}})^n = \mathfrak{m}^n\) . \(\hat{\mathbf{A}}\) 对每个整数 \(n > 0\) 都成立，且 \(\hat{\mathbf{A}}\) 上的拓扑是 m-进拓扑。

记 \(A_{n} = m^{n}\)，它是 A 的有限生成理想。公式 \(m^{p}A_{n} = m^{n+p}\) 表明，\(A_{n}\) 上由 m-进拓扑诱导的拓扑，与 A-模 \(A_{n}\) 上的 m-进拓扑相同（第 1 节，例 3）。将命题 16 应用于 \(E = A_{n}\)，得到 \(\hat{A}_{n} = \hat{A}\) . \(A_{n}\)，换言之，\(\widehat{m^{n}} = m^{n}\) . \(\hat{A}\)。特别地，\(\hat{m} = m\) . \(\hat{A}\)，从而

\[
(\hat {\mathfrak {m}}) ^ {n} = \mathfrak {m} ^ {n}. \hat {\mathbf {A}} = \hat {\mathbf {A}}. \mathbf {A} _ {n}
\]

（参见~练习 12）。

\textbf{滤过环的 Hausdorff 完备化的例子}\par

\begin{enumerate}
\def\labelenumi{(\arabic{enumi})}
\tightlist
\item
  设 \(\mathbf{A}\) 是一个 \(\mathbf{N}\) 型分次环，且 \((\mathbf{A}_n)_{n\geqslant 0}\) 是其分次；给它赋予相伴滤过，该滤过是分离且穷竭的（第 1 节，例 1）。加法群 \(\mathbf{A}\) 被典范地识别为 \(\mathbf{B}_{n}\) 的一个子群，仿佛 \(n\) 是由该子群定义的

\(B = \prod_{n \in N} A_n\)；若在 \(B\) 上赋予离散拓扑的乘积拓扑，则 \(A\) 上诱导的拓扑就是由 \(A\) 上的滤过定义的拓扑；此外，\(B\) 是完备拓扑群，而 \(A\) 在 \(B\) 中稠密（《一般拓扑学》，第三章 § 2，第 9 节，命题 25）。于是加法拓扑群 \(B\) 可识别为 \(\hat{A}\)（即 Hausdorff 加法群 \(A\) 的完备化），由命题 15 可知，它具有唯一的环结构，使其成为拓扑环 \(A\) 的完备化。为定义这个环中的乘法，

注意，若记 \(A_{n}^{\prime} = \sum_{i > n} A_{i}\)，则两侧理想 \(A_{n}^{\prime}\) 在 B 中的闭包是集合 \(B_{n}\)，其中的元素 \(x = (x_{i}) \in \mathrm{B}\) 满足 \(x_{i} = 0\)（对 \(i \leqslant n\)）。现在令 \(x = (x_{i}), y = (y_{i})\) 为 B 的两个元素，令 \(z = (z_{i})\) 为它们的乘积。那么，对所有 n \textgreater{} 0，有 \(x \equiv x_{n}^{\prime} \pmod{B_{n}}\)、\(y \equiv y_{n}^{\prime} \pmod{B_{n}}\)，其中 \(x_{n}^{\prime} = (x_{i})_{0 \leqslant i \leqslant n}, y_{n}^{\prime} = (y_{i})_{0 \leqslant i \leqslant n}\)，从而 \(z \equiv x_{n}^{\prime} y_{n}^{\prime} \pmod{B_{n}}\)。但 \(x_{n}^{\prime}\) 和 \(y_{n}^{\prime}\) 属于 A，因而可见，对所有 \(n \in N\)，

\[
z _ {n} = \sum_ {j = 0} ^ {n} x _ {j} y _ {n - j}.\tag{26}
\]

特别地，我们再次得到第 6 节命题 6 的推论：如果 C 是交换环，则带有相应滤过的多项式环 \(C[X_{1},\ldots,X_{r}]\) 的完备化（按总次数的通常分次）可典范地识别为形式幂级数环 \(C[[X_{1},\ldots,X_{r}]]\)（参见~《代数》，第四章 § 5，第 10 节）。
\end{enumerate}

* (2) 设 K 是一个带赋值的完备交换域。K 上 r 个变量的收敛级数环的完备化可典范地识别为形式幂级数环 \(K[[X_{1},\ldots,X_{r}]]\)。 *

\begin{enumerate}
\def\labelenumi{(\arabic{enumi})}
\setcounter{enumi}{2}
\tightlist
\item
  设 \(\alpha\) 是主理想整环中的非零非可逆元；A 上的 \((\alpha)\)-进拓扑也称为 \(\alpha\)-进拓扑；它是 Hausdorff 的，因为理想 \((\alpha^n)\) 的交为 0（《代数》，第七章 §1，第 3 节）。注意，A 关于此拓扑的完备化不一定是整环（参见~第 13 节，注 3）。相伴分次环 \(\mathrm{gr}(\mathbf{A}) = \mathrm{gr}(\hat{\mathbf{A}})\) 可典范地同构于 \((\mathbf{A}/\mathfrak{a})[\mathbf{X}]\)（第 3 节，例 1）。若 \(\mathbf{A} = \mathbf{Z}\)，则 \(\mathbf{Z}\) 关于 \(n\)-进拓扑（\(n > 1\)）的完备化记为 \(\mathbf{Z}_n\)，其元素称为 \(n\)-进整数。
\end{enumerate}

每个 \(\mathbf{Z} / n^k\mathbf{Z}\) 的元素都有唯一的形如 \(\sum_{i=0}^{k-1} a_i n^i\) 的代表元，其中 \(0 \leqslant a_i \leqslant n - 1\) 对所有 \(i\) 成立；此外，它在 \(\mathbf{Z} / n^{k-1}\mathbf{Z}\) 中的典范像是 \(\sum_{i=0}^{k-2} a_i n^i\) 的类。这些说明以及 \(\mathbf{Z}_n\) 可典范地识别为逆极限 \(\lim_{k} \mathbf{Z}/n^k \mathbf{Z}\) 这一事实，立即表明每个

\(\mathbf{Z}_n\) 中的元素都能唯一写成 \(\sum_{i=0}^{\infty} a_i n^i\) 的形式，其中 \(0 \leqslant a_i < n\)，反之这样的级数在 \(\mathbf{Z}_n\) 中收敛。
% semantic-fragments: legacy-input-seam
\relax{}
